\section*{Capítulo \ref{ch:b1:s-block} --- El hidrógeno y el bloque s}

\begin{solution}{exo:b1:s-block:1}
\ce{KH} y \ce{CaH2}: salinos; \ce{SiH4} y \ce{H2S}: covalentes;
\ce{PdH_{0.6}}: metálico. \ce{CaH2 + 2H2O -> Ca(OH)2 + 2H2}.
\end{solution}

\begin{solution}{exo:b1:s-block:2}
\ce{2Li + 2H2O -> 2LiOH + H2}, e igual con Na y K;
\ce{Mg + 2HCl -> MgCl2 + H2}.
\end{solution}

\begin{solution}{exo:b1:s-block:3}
Básicos: \ce{Na2O}, \ce{MgO} (\ce{MgO + 2HCl -> MgCl2 + H2O}). Ácidos:
\ce{SiO2}, \ce{SO3}, \ce{CO2} (\ce{SO3 + H2O -> H2SO4}). Anfóteros:
\ce{Al2O3}, \ce{ZnO} (\ce{ZnO + 2OH- + H2O -> Zn(OH)4^2-}).
\end{solution}

\begin{solution}{exo:b1:s-block:4}
A pH 14, el par del agua es \ce{2H2O + 2e- -> H2 + 2OH-}, $E = -0.83$~V;
$\log K = 2(-0.83 + 2.71)/0.059 = 64$.
\end{solution}

\begin{solution}{exo:b1:s-block:5}
En la llama, los choques llevan un electrón a un nivel excitado; al
volver al \omterm{def:b1:quantum-numbers:energy-level}{estado fundamental}, el átomo emite un fotón de energía igual a
la diferencia, propia del elemento. $E = hc/\lambda = 6.626 \times 10^{-34} \times
3.00 \times 10^8/(589 \times 10^{-9}) = \qty{3.375e-19}{J} = \qty{2.11}{eV}$.
% ledger: const:h, const:c, const:eV
\end{solution}

\begin{solution}{exo:b1:s-block:6}
$10^6/106.0 = \qty{9.43e3}{mol}$ de carbonato de sodio: \qty{9.43e3}{mol}
de caliza, \qty{0.944}{t}, y $2 \times 9.43 \times 10^3/0.90 =
\qty{2.10e4}{mol}$ de sal, \qty{1.23}{t}.
\end{solution}

\begin{solution}{exo:b1:s-block:7}
$[\ce{Mg^2+}] = 1.3/24.3 = \qty{0.053}{mol/L}$; $\mathrm{pH} = 14.00 -
\frac12(11.25 + \log 0.053) = 9.0$. Los \qty{53.5}{mol} de magnesio de
\qty{1.0}{m^3} necesitan \qty{53.5}{mol} de \ce{Ca(OH)2}: \qty{4.0}{kg}.
\end{solution}

\begin{solution}{exo:b1:s-block:8}
La constante de equilibrio (termodinámica) es enorme en ambos casos; la
velocidad (cinética) no la decide $E^\circ$. El litio funde a una
temperatura más alta, sigue sólido y se recubre de una capa de producto
menos soluble, y pierde su electrón con menos facilidad que el potasio
(energía de ionización más alta): reacciona de forma regular, y el
potasio, violentamente.
\end{solution}

\begin{solution}{exo:b1:s-block:9}
$10^6/100.1 = \qty{9.99e3}{mol}$: CaO, $9.99 \times 10^3 \times 56.1 =
\qty{0.560}{t}$; \ce{CO2}, \qty{9.99e3}{mol}, $V = nRT/p = 9.99 \times 10^3
\times 8.314 \times 298.15/10^5 = \qty{248}{m^3}$.
\end{solution}

\begin{solution}{exo:b1:s-block:10}
$2.9 \times 10^8/6.9 = \qty{4.2e7}{kmol}$ de Li (masa en kg), $\times 73.8/2 =
\qty{1.55e9}{kg}$: unos 1.6 millones de toneladas de carbonato de litio.
$2.9 \times 10^8/8 = \num{3.6e7}$ baterías.
\end{solution}

\begin{solution}{exo:b1:s-block:11}
1.31, 0.84 y 0.71 por ciento a 20, 80 y \qty{100}{\celsius}: la
\omterm{def:b1:precipitation:solubility}{solubilidad} disminuye. En caliente queda menos carbonato en disolución:
se recupera más.
\end{solution}

\begin{solution}{exo:b1:s-block:12}
El \ce{NaHCO3} es la menos soluble de las sales que pueden formar los
iones, y la disolución se satura de ella en primer lugar. El amoniaco
hace la disolución lo bastante básica para convertir el \ce{CO2} en
\ce{HCO3-} y, a la vez, puede recuperarse (como \ce{NH4Cl}, que la cal
regenera); el hidróxido de sodio se consumiría y costaría más que el
producto.
\end{solution}

\begin{solution}{pb:b1:s-block:1}
\textbf{1.} Li, \qty{6.0}{g/L}; Mg, \qty{2.4}{g/L}.
\textbf{2.} Cuatro metros cúbicos de salmuera dan uno: se evaporan tres
metros cúbicos de agua.
\textbf{3.} El sol y el clima seco aportan la energía gratis.
\textbf{4.} El cloruro de sodio, con mucho la sal más abundante, se
satura primero.
\textbf{5.} Li, $6000/6.9 = \qty{870}{mol}$; Mg, $2400/24.3 = \qty{98.8}{mol}$.
\textbf{6.} El magnesio precipitaría con el carbonato y contaminaría el
producto.
\textbf{7.} \ce{Mg^2+ + Ca(OH)2 -> Mg(OH)2 + Ca^2+}.
\textbf{8.} $98.8 \times 74.1 = \qty{7.32}{kg}$.
\textbf{9.} Quedan $[\ce{Mg^2+}] = \num{9.88e-5}$ mol/L: $\mathrm{pH} = 14.00 -
\frac12(11.25 - 4.01) = 10.38$.
\textbf{10.} El hidróxido de litio es soluble: a este pH no se forma
ningún sólido.
\textbf{11.} \ce{Ca^2+ + CO3^2- -> CaCO3}.
\textbf{12.} $98.8 \times 106.0 = \qty{10.5}{kg}$.
\textbf{13.} \ce{2Li+ + CO3^2- -> Li2CO3}.
\textbf{14.} $435 \times 106.0 = \qty{46.1}{kg}$.
\textbf{15.} $435 \times 73.8 = \qty{32.1}{kg}$.
\textbf{16.} El carbonato de litio es menos soluble en caliente.
\textbf{17.} $0.0084 \times 1000 = \qty{8.4}{kg}$.
\textbf{18.} Con pequeños volúmenes de agua caliente, en la que es menos
soluble.
\textbf{19.} $32.1 - 8.4 = \qty{23.7}{kg}$.
\textbf{20.} $23.7/32.1 = \qty{74}{\percent}$.
\textbf{21.} Todavía contiene litio: se devuelve a las piscinas.
\textbf{22.} $2.9 \times 10^8/6.0 = \num{4.8e7}$ metros cúbicos.
\textbf{23.} \textbf{\qty{24}{kg}} (23.7) de carbonato de litio por metro
cúbico.
\end{solution}
